% !PRINTED 10
\weberpage{10}

%<<CONT
that $\alpha$ be the smallest positive exponent for which $a^\alpha = 1$.
We have, always and only then,
\[
a^\mu = a^{\mu'},
\]
when $\mu \equiv \mu'\ (\mathrm{mod}\ \alpha)$, and here the exponents may
also be negative, if $a^{-r}$ denotes the $r$th power of the element
$a^{-1}$, or the element opposite to $a^r$.

The sequence
\[
A = 1,\ a,\ a^2 \ldots a^{\alpha-1}
\]
then contains $\alpha$ elements of $P$ different from one another, among
which all powers of $a$ are represented. But the elements $A$ evidently
form a group of degree $\alpha$, because the exponents simply add under
composition. This group is a divisor of $P$, and hence $\alpha$ is
\gesperrt{a divisor of $n$}. Hence for every element $a^n = 1$.

\gesperrt{The group $A$ is called the period of the element $a$, and $\alpha$ is
also called the degree of the element $a$.}

\begin{center}
\S.\,3.

\gesperrt{Normal divisors of a group.}
\end{center}

If $P$ is, as above, a group and $Q$ a divisor of $P$, and further $b$ an
element of $P$ not contained in $Q$, then the subsidiary group $Qb$ is not
a group. On the other hand the system $b^{-1}Qb$, or written out at
length
\[
1,\ b^{-1}a_1 b,\ b^{-1}a_2 b \ldots b^{-1}a_{\nu-1} b
\]
certainly forms a group, because
\[
b^{-1}a_1 b\ .\ b^{-1}a_2 b = b^{-1}a_1 a_2 b
\]
and this group is isomorphic with $Q$. The group $b^{-1}Qb$ is called
\gesperrt{the group transformed by $b$}. If $b$ itself belongs to $Q$, then
$b^{-1}Qb$ is identical with $Q$. If one takes for $b$ the various
elements of $P$, one obtains a whole family of such groups, which we call
the \gesperrt{divisors of $P$ conjugate to $Q$}, or also briefly \gesperrt{conjugate groups}.

It may happen that all the conjugate divisors are identical with one
another. We have already seen in the case of the permutation groups that
this case is of particular importance, and we therefore now introduce in
general the following definition:
%CONT>>
